% Part: applied-modal-logics
% Chapter: epistemic-logic
% Section: public-announcement-logic-semantics

\documentclass[../../../include/open-logic-section]{subfiles}

\begin{document}

\olfileid[pt]{aml}{el}{psm}

\olsection{Semântica da lógica de anúncios públicos}

Os modelos relacionais das lógicas de anúncios públicos são os mesmos
das lógicas epistêmicas. A semântica do operador de anúncio público,
porém, é nova.

\begin{defn}\ollabel{defn:mmodels} A \emph{verdade de !!a{formula}~$!A$
  em~$w$} em um~$\mModel M = \tuple{W, R, V}$, em símbolos
  $\mSat{M}{!A}[w]$, é definida indutivamente da seguinte maneira:
  \begin{enumerate}
  \tagitem{prvFalse}{\indcase{!A}{\lfalse}{Nunca $\mSat{M}{\lfalse}[w]$}.}{}
  \tagitem{prvTrue}{\indcase{!A}{\ltrue}{Sempre $\mSat{M}{\ltrue}[w]$}.}{}
  \item $\mSat{M}{p}[w]$ sse $w \in V(p)$
  \tagitem{prvNot}{\indcase{!A}{\lnot !B}{$\mSat{M}{\indfrm}[w]$ sse
    $\mSat/{M}{!B}[w]$}.}{}
  \tagitem{prvAnd}{\indcase{!A}{(!B \land !C)}{$\mSat{M}{\indfrm}[w]$ sse
    $\mSat{M}{!B}[w]$ e $\mSat{M}{!C}[w]$}.}{}
  \tagitem{prvOr}{\indcase{!A}{(!B \lor !C)}{$\mSat{M}{\indfrm}[w]$ sse
    $\mSat{M}{!B}[w]$ ou $\mSat{M}{!C}[w]$} (ou ambos).}{}
  \tagitem{prvIf}{\indcase{!A}{(!B \lif !C)}{$\mSat{M}{\indfrm}[w]$ sse
    $\mSat/{M}{!B}[w]$ ou $\mSat{M}{!C}[w]$}.}{}
  \tagitem{prvIff}{\indcase{!A}{(!B \liff !C)}{$\mSat{M}{\indfrm}[w]$ sse
    ou ambos $\mSat{M}{!B}[w]$ e $\mSat{M}{!C}[w]$, ou
    nem $\mSat{M}{!B}[w]$ nem $\mSat{M}{!C}[w]$}.}{}
  \item\ollabel{defn:sub:mmodels-box}
    \indcase{!A}{\Knows_a !B}{$\mSat{M}{\indfrm}[w]$ sse
    $\mSat{M}{!B}[w']$ para todo $w' \in W$ tal que $R_a ww'$}
  \item\ollabel{defn:sub:mmodels-pal}
    \indcase{!A}{[!B] !C}{$\mSat{M}{\indfrm}[w]$ sse
    $\mSat{M}{!B}[w]$ implica $\mSat{M \mid !B}{!C}[w]$}

  
  Em que $\mModel M \mid !B = \tuple{W', R', V'}$ é definido da
  seguinte maneira:
  \begin{enumerate}
    \item $W' = \Setabs{ u \in W}{\mSat{M}{!B}[u]}$. Assim, os mundos de
      $\mModel M \mid !B$ são os mundos de~$\mModel M$ nos quais $!B$
      vale.
    \item $R'_a = R_a \cap (W' \times W')$. A relação de acessibilidade
      de cada agente é simplesmente restringida aos mundos que permanecem em~$W'$.
    \item $V'(p) = \Setabs{ u \in W'}{u \in V(p) }$. De modo semelhante,
      as valorações proposicionais nos mundos permanecem as mesmas, o que
      representa a ideia de que eventos informacionais não alteram o
      valor de verdade das variáveis proposicionais.
  \end{enumerate}
  
  \end{enumerate} 
\end{defn}

O aspecto distintivo das lógicas de anúncios públicos, portanto, é que
a verdade de !!a{formula} em~$\mModel M$ às vezes só pode ser decidida
por referência a um modelo diferente do próprio~$\mModel M$.

Observe também que nossa semântica trata o operador de anúncio como um
operador~$\Box$; assim, se !!a{formula}~$!A$ não pode ser anunciada de
modo verídico em um mundo, então $[!A]B$ valerá ali trivialmente, assim
como todas as !!{formula}s $\Box$ valem em pontos terminais. 

\begin{figure}
  \begin{center}
    \begin{tikzpicture}[modal]
      \node[world] (w1) [label={below left:$p, \lnot q$}]{$w_1$}; 
      \node[world] (w2) [left=of w1, label={left:$\lnot p, \lnot q$}]{$w_2$}; 
      \node[world] (w3) [above=of w1, label={right:$p, q$}] {$w_3$};
      \draw[<->] (w1) to [swap] node {$b$} (w2);
      \draw[->,loop below] (w1) to node {$a, b$} (w1);
      \draw[<->] (w1) to [swap] node {$a$} (w3);
      \draw[->,loop above] (w2) to node {$a, b$} (w2) ;
      \draw[->,loop above] (w3) to node {$a, b$} (w3) ;
      
      \node[world](v1)[right of=w1, xshift=1.25in, label={right:$p, \lnot q$}]{$w'_1$};
      \node[world](v3)[above of=v1,label={right:$p,q$}]{$w'_3$};
      \draw[<->] (v1) to node {$a$} (v3);
      \draw[->,loop below] (v1) to node {$a,b$} (v1);
      \draw[->,loop above] (v3) to node {$a,b$} (v3) ;
      
      \node(m)[below=of w1]{$\mModel M$};
      \node(m')[below=of v1]{$\mModel M \mid p$}; 

      \draw[-,dotted] (w1) to node {\footnotesize anúncio de $p$}(v1) ;
     
    \end{tikzpicture}
  \end{center}
  \caption{Antes e depois do anúncio público de $p$.}
  \ollabel{fig:announcement-example}
\end{figure}

Podemos entender o anúncio público de !!a{formula} como uma redução do
modelo, ou como sua restrição aos mundos nos quais a !!{formula} era
verdadeira. \olref{fig:announcement-example} apresenta um exemplo dos
efeitos do anúncio público de~$p$. Um aspecto notável desse modelo é
que o agente~$b$ passa a saber que~$p$ como resultado do anúncio,
enquanto o agente~$a$ não (pois $a$ já sabia que~$p$ era verdadeira).

Mais formalmente, temos $\mSat{M}{\lnot \Knows_b p}[w_1]$, mas $\mSat{M
\mid p}{\Knows_b p}[w'_1]$. Isso implica $\mSat{M}{[p] \Knows_b
p}[w_1]$. Podemos, porém, fazer afirmações ainda mais fortes sobre o
resultado do anúncio. De fato,
$\mSat{M}{[p]\CKnows_{\{a,b\}} p}[w_1]$. Em outras palavras, depois que
$p$ é anunciada, ela se torna \emph{conhecimento comum}.

Podemos perguntar, porém, se isso vale no caso geral e se um anúncio
verídico de~$!A$ fará \emph{sempre} com que $!A$ se torne conhecimento
comum. Talvez seja surpreendente que a resposta seja não. De fato, é
possível anunciar de modo verídico !!{formula}s que deixam de ser
verdadeiras assim que são anunciadas. Considere, por exemplo, os efeitos
de anunciar $p \land \lnot \Knows_b p$ em~$w_1$ na
\olref{fig:announcement-example}. Na realidade, $\mModel M \mid p$ e
$\mModel M \mid (p \land \lnot \Knows_b p)$ são o mesmo modelo.
Entretanto, como já observamos, $\mSat{M \mid p}{\Knows_b
p}[w'_1]$. Portanto, $\mSat{M \mid (p \land \lnot \Knows_b p)}{\lnot
(p \land \lnot \Knows_b p)}[w'_1]$; assim, essa é !!a{formula} que se
torna falsa depois de ser anunciada. 

\end{document}
